\documentclass[10pt]{beamer}

%------------------------------------------------
% Tema
%------------------------------------------------
\usetheme{moloch}
\definecolor{mPurple}{RGB}{91,48,105}
\setbeamercolor{palette primary}{bg=mPurple}
\setbeamercolor{normal text}{fg=black}

%------------------------------------------------
% Paquetes esenciales
%------------------------------------------------
\usepackage[utf8]{inputenc}
\usepackage[T1]{fontenc}
\usepackage[spanish]{babel}

\usepackage{amsmath, amsfonts, amssymb}
\usepackage{graphicx}
\usepackage{booktabs}

% Código
\usepackage{listings}
\usepackage{xcolor}

%------------------------------------------------
% Configuración de código (R)
%------------------------------------------------
\lstset{
  language=R,
  basicstyle=\ttfamily\small,
  keywordstyle=\color{blue},
  commentstyle=\color{gray},
  stringstyle=\color{red},
  numbers=left,
  numberstyle=\tiny,
  stepnumber=1,
  numbersep=5pt,
  backgroundcolor=\color{gray!10},
  showspaces=false,
  showstringspaces=false,
  showtabs=false,
  frame=single,
  breaklines=true,
  breakatwhitespace=true
}

%------------------------------------------------
% Información del documento
%------------------------------------------------
\title{Pre-Curso de R para Stress Testing en Pensiones}
\author{MACROPOL}
\date{\today}

%------------------------------------------------
\begin{document}

%------------------------------------------------
\begin{frame}
\titlepage
\end{frame}

%------------------------------------------------
%\begin{frame}{Contenido}
%\tableofcontents
%\end{frame}
%------------------------------------------------

%------------------------------------------------
\subsection{Objetivo}

\begin{frame}{Pre-Curso de R para Stress Testing}

\textbf{Objetivo:}
\begin{itemize}
\item Preparar a los participantes para implementar ejercicios en R
\item Enfocado en simulación y análisis de riesgo
\end{itemize}

\vspace{0.5cm}

\textbf{Importante:}
\begin{itemize}
\item No es un curso de programación
\item Es un curso de \textbf{R aplicado a pensiones}
\end{itemize}

\end{frame}

%------------------------------------------------
\begin{frame}{¿Qué necesitamos de R?}

\textbf{Capacidad mínima:}

\begin{itemize}
\item Ejecutar código
\item Modificar parámetros
\item Interpretar resultados
\end{itemize}

\vspace{0.5cm}

\textbf{Aplicación en el curso:}
\begin{itemize}
\item Simulación Monte Carlo
\item Cálculo de métricas (VaR, ES)
\item Construcción de escenarios
\end{itemize}

\end{frame}

%------------------------------------------------
\begin{frame}[t]{Agenda}

\begin{enumerate}
\item Instalación
\item Variables y vectores
\item Estructuras de control
\item Funciones
\item Simulación aleatoria
\item Métricas de riesgo
\item Visualización
\end{enumerate}

\end{frame}

%------------------------------------------------
%------------------------------------------------
\begin{frame}{Hilo conductor del pre-curso}

\textbf{Objetivo práctico}

\vspace{0.3cm}
Construiremos paso a paso un modelo de riesgo previsional:

\begin{enumerate}
\item Ahorro individual (determinístico)
\item Incorporar incertidumbre (retornos)
\item Simular múltiples escenarios
\item Medir riesgo (VaR)
\item Visualizar distribución
\end{enumerate}

\vspace{0.5cm}

\textbf{Resultado final:}
Capacidad de hacer un stress testing básico en pensiones

\end{frame}

%---------------------------------------------



\begin{frame}{Instalación de R y RStudio}

\textbf{Contexto}

\vspace{0.2cm}
Antes de implementar modelos de stress testing, es necesario contar con un entorno reproducible en R.

\vspace{0.4cm}

\textbf{Herramientas}
\begin{itemize}
\item R: Motor de cálculo estadístico
\item RStudio: Entorno de desarrollo (IDE)
\item Paquetes clave: \texttt{tidyverse, quantmod, forecast, vars, rugarch}
\end{itemize}

\vspace{0.4cm}

\textbf{Instalación}
\begin{itemize}
\item Descargar R: \texttt{https://cran.r-project.org}
\item Descargar RStudio: \texttt{https://posit.co/download/rstudio-desktop/}
\item Instalar con configuración por defecto
\end{itemize}

\end{frame}

%------------------------------------------------

\begin{frame}[fragile]{Configuración inicial en R}
\footnotesize
\textbf{Verificación}

\begin{lstlisting}
version
\end{lstlisting}

\vspace{0.3cm}

\textbf{Instalación de paquetes}

\begin{lstlisting}
install.packages(c(
 "tidyverse",
 "quantmod",
 "forecast",
 "vars",
 "rugarch",
 "PerformanceAnalytics"
))
\end{lstlisting}


\textbf{Carga de librerías}

\begin{lstlisting}
library(tidyverse)
library(quantmod)
library(forecast)
library(vars)
library(rugarch)
library(PerformanceAnalytics)
\end{lstlisting}

\end{frame}

%------------------------------------------------
\begin{frame}[fragile]{1. Variables y Vectores}

\textbf{Asignación:}

\begin{verbatim}
x <- 10
\end{verbatim}

\vspace{0.3cm}

\textbf{Vectores:}

\begin{verbatim}
v <- c(1,2,3,4,5)
\end{verbatim}

\vspace{0.3cm}

\textbf{Operaciones:}

\begin{verbatim}
mean(v) # media
sum(v)  # suma
\end{verbatim}

\vspace{0.3cm}


\textbf{Ejercicio 1}
    \begin{itemize}
        \item Crear un vector 
        \item Calcular su media
    \end{itemize}
\end{frame}

%------------------------------------------------
\begin{frame}[fragile]{2. Estructuras de Control}

\textbf{Loop:}

\begin{verbatim}
for(i in 1:10){
  print(i)
}
\end{verbatim}

\vspace{0.3cm}

\textbf{Condicional:}

\begin{verbatim}
if(x > 0){
  print("positivo")
}
\end{verbatim}

\vspace{0.3cm}

\textbf{Aplicación:}
Simulación de trayectorias

\end{frame}

%------------------------------------------------
\begin{frame}{Ejercicio 2}

\textbf{Tarea:}

\begin{itemize}
\item Simular la evolución de un ahorro durante 10 años
\end{itemize}

\vspace{0.3cm}

\textbf{Pregunta:}
¿Qué representa cada iteración del loop?

\end{frame}

%------------------------------------------------
\begin{frame}[fragile]{3. Funciones}

\textbf{Definición:}

\begin{verbatim}
simulate <- function(W){
  return(W * 1.05)
}
\end{verbatim}

\vspace{0.3cm}

\textbf{Uso:}

\begin{verbatim}
simulate(1000)
\end{verbatim}

\textbf{Interpretación:}

\begin{itemize}
\item Cada iteración representa un período (mes/año)
\item Estamos construyendo acumulación de riqueza
\end{itemize}

\vspace{0.3cm}

\textbf{Siguiente paso:}
Formalizar este proceso en una función reutilizable
\end{frame}
%------------------------------------------------
%------------------------------------------------
\begin{frame}{De código a modelo previsional}

\textbf{Intuición}

\vspace{0.3cm}
Hasta ahora hemos visto:

\begin{itemize}
\item Variables
\item Loops
\item Funciones
\end{itemize}

\vspace{0.4cm}

\textbf{Ahora:}

Usaremos estas herramientas para construir un modelo de acumulación de pensiones

\end{frame}

%------------------------------------------------
\begin{frame}{Ejercicio 3: Función de riqueza final}

\textbf{Contexto}

\vspace{0.2cm}
En el sistema de capitalización individual, la riqueza acumulada depende de:
\begin{itemize}
\item Aportes periódicos
\item Rentabilidad del portafolio
\item Horizonte de inversión
\end{itemize}

Este tipo de dinámica es la base de cualquier stress testing previsional.

\vspace{0.4cm}

\textbf{Modelo}

La riqueza final puede aproximarse como:

\begin{itemize}
\item Acumulación de aportes capitalizados
\item Bajo una tasa de retorno constante (baseline)
\end{itemize}

\vspace{0.3cm}

\textbf{Tarea (Paso a paso)}

\begin{enumerate}
\item Definir una función en R que reciba:
\begin{itemize}
\item aporte mensual (\texttt{c})
\item tasa de retorno (\texttt{r})
\item número de períodos (\texttt{n})
\end{itemize}

\item Calcular la riqueza acumulada iterativamente
\end{enumerate}

\end{frame}

%------------------------------------------------

\begin{frame}[fragile]{Implementación en R}

\textbf{Paso 1: Definir la función}

\begin{lstlisting}
riqueza_final <- function(c, r, n) {
  W <- 0
  for (t in 1:n) {
    W <- (W + c) * (1 + r)
  }
  return(W)
}
\end{lstlisting}

\vspace{0.4cm}

\textbf{Paso 2: Ejecutar el modelo (baseline)}

\begin{lstlisting}
riqueza_final(c = 100, r = 0.005, n = 360)
\end{lstlisting}

\vspace{0.3cm}

\textbf{Interpretación}
\begin{itemize}
\item Representa acumulación en escenario base (sin estrés)
\item Supone retorno constante
\end{itemize}

\end{frame}

%------------------------------------------------
%------------------------------------------------
\begin{frame}{Limitación del modelo baseline}

\textbf{Problema}

\vspace{0.3cm}
El modelo anterior asume:

\begin{itemize}
\item Retornos constantes
\item No hay incertidumbre
\end{itemize}

\vspace{0.4cm}

\textbf{Realidad}

\begin{itemize}
\item Los mercados son volátiles
\item Los retornos son inciertos
\end{itemize}

\vspace{0.4cm}

\textbf{Solución:}
Introducir simulación aleatoria

\end{frame}
%------------------------------------------------
\begin{frame}[fragile]{4. Simulación Aleatoria}

\textbf{Normal:}

\begin{verbatim}
rnorm(100, mean=0.05, sd=0.15)
\end{verbatim}

\vspace{0.3cm}

\textbf{Interpretación:}
Cada número representa un posible retorno

\vspace{0.3cm}

\textbf{Aplicación:}
Monte Carlo

\end{frame}
%-------------------------------------------------

\begin{frame}[fragile]{Ejercicio 4: Escenario adverso}

\textbf{Narrativa económica}

\vspace{0.2cm}
Shock financiero:
\begin{itemize}
\item Caída en retornos
\item Mayor volatilidad
\end{itemize}

\textbf{Modificación del modelo}

\begin{lstlisting}
set.seed(123)

riqueza_estocastica <- function(c, mu, sigma, n) {
  W <- 0
  for (t in 1:n) {
    r_t <- rnorm(1, mean = mu, sd = sigma)
    W <- (W + c) * (1 + r_t)
  }
  return(W)
}
\end{lstlisting}

\vspace{0.3cm}

\textbf{Ejecutar escenario adverso}

\begin{lstlisting}
riqueza_estocastica(100, 0.002, 0.02, 360)
\end{lstlisting}

\end{frame}

%------------------------------------------------

\begin{frame}{Discusión conceptual}

\textbf{Pregunta clave}

\vspace{0.3cm}
¿Por qué es útil encapsular el modelo en una función?

\vspace{0.4cm}

\textbf{Respuesta esperada}

\begin{itemize}
\item Permite reutilizar el modelo en múltiples escenarios
\item Facilita simulaciones Monte Carlo
\item Mejora trazabilidad y auditoría (clave regulatoria)
\item Separa supuestos (inputs) de resultados (outputs)
\end{itemize}

\vspace{0.4cm}

\textbf{Conexión con stress testing}

\begin{itemize}
\item Baseline vs adverso se define cambiando parámetros
\item Permite escalar a miles de simulaciones
\item Base para análisis de distribución de pensiones futuras
\end{itemize}

\textbf{Limitación}

\begin{itemize}
\item Una sola simulación no es suficiente
\item Necesitamos una distribución de resultados
\end{itemize}

\vspace{0.3cm}

\textbf{Siguiente paso:}
Simular muchas trayectorias (Monte Carlo)

\end{frame}
%------------------------------------------------
\begin{frame}[fragile]{5. Métricas de Riesgo}

\textbf{Cálculo:}

\begin{verbatim}
mean(x)
quantile(x, 0.05)
\end{verbatim}

\vspace{0.3cm}

\textbf{Conceptos:}

\begin{itemize}
\item Media: resultado esperado
\item VaR: riesgo
\end{itemize}

\end{frame}

%------------------------------------------------
\begin{frame}{Ejercicio 5: VaR sobre riqueza previsional}

\textbf{Contexto}

\vspace{0.2cm}
A partir del ejercicio anterior, ahora contamos con un modelo que genera riqueza bajo incertidumbre.

\vspace{0.2cm}
El siguiente paso en stress testing es medir el riesgo de resultados adversos.

\vspace{0.4cm}

\textbf{Tarea}

\begin{enumerate}
\item Simular múltiples trayectorias de riqueza final
\item Construir la distribución de resultados
\item Calcular el Value at Risk (VaR)
\end{enumerate}

\end{frame}

%------------------------------------------------

\begin{frame}[fragile]{Implementación en R}

\textbf{Paso 1: Simulación Monte Carlo}

\begin{lstlisting}
set.seed(123)

N <- 1000

simulaciones <- replicate(N,
  riqueza_estocastica(100, 0.002, 0.02, 360)
)
\end{lstlisting}

\vspace{0.4cm}

\textbf{Paso 2: Calcular VaR (percentil 5\%)}

\begin{lstlisting}
VaR_95 <- quantile(simulaciones, 0.05)
VaR_95
\end{lstlisting}

\vspace{0.3cm}

\textbf{Interpretación}
\begin{itemize}
\item Representa el nivel de riqueza en un escenario adverso extremo
\item 5\% de los casos tienen resultados peores
\end{itemize}

\end{frame}

%------------------------------------------------

\begin{frame}{Discusión: Relevancia del VaR en pensiones}

\textbf{Pregunta}

\vspace{0.3cm}
¿Por qué el VaR es relevante en pensiones?

\vspace{0.4cm}

\textbf{Respuesta esperada}

\begin{itemize}
\item Mide riesgo de insuficiencia de ahorro previsional
\item Permite evaluar escenarios adversos (cola de la distribución)
\item Es clave para stress testing regulatorio
\item Ayuda a identificar vulnerabilidad de afiliados
\end{itemize}

\vspace{0.4cm}

\textbf{Limitación importante}

\begin{itemize}
\item No captura severidad más allá del percentil (tail risk)
\item Debe complementarse con Expected Shortfall
\end{itemize}

\end{frame}

%------------------------------------------------

\begin{frame}{Interpretación para política pública}

\textbf{Intuición económica}

\begin{itemize}
\item El VaR traduce incertidumbre financiera en riesgo social
\item Permite cuantificar probabilidad de pensiones insuficientes
\end{itemize}

\vspace{0.4cm}

\textbf{Aplicación regulatoria}

\begin{itemize}
\item Definición de límites de riesgo en portafolios previsionales
\item Evaluación de sostenibilidad del sistema
\item Diseño de políticas de mitigación (garantías mínimas)
\end{itemize}

\vspace{0.4cm}

\textbf{Conexión con el curso}

\begin{itemize}
\item Módulo 3: Simulación de escenarios
\item Módulo 4: Riesgo de portafolio
\item Módulo 6: Reporte regulatorio
\end{itemize}

\textbf{Limitación del VaR}

\begin{itemize}
\item Es una medida puntual
\item No muestra la forma completa del riesgo
\end{itemize}

\vspace{0.3cm}

\textbf{Siguiente paso:}
Visualizar toda la distribución
\end{frame}
%------------------------------------------------
\begin{frame}[fragile]{6. Visualización}

\textbf{Histograma:}

\begin{verbatim}
hist(x)
\end{verbatim}

\vspace{0.3cm}

\textbf{Aplicación:}
Distribución de resultados

\end{frame}

%------------------------------------------------
\begin{frame}{Ejercicio 6: Distribución de retornos}

\textbf{Contexto}

\vspace{0.2cm}
Luego de simular múltiples trayectorias, es fundamental visualizar la distribución de resultados para identificar riesgos.

\vspace{0.4cm}

\textbf{Tarea}

\begin{enumerate}
\item Obtener retornos simulados
\item Graficar su distribución
\item Identificar zonas de riesgo
\end{enumerate}

\end{frame}

%------------------------------------------------

\begin{frame}[fragile]{Implementación en R}

\textbf{Paso 1: Usar simulaciones anteriores}

\begin{lstlisting}
hist(simulaciones, breaks = 30,
     main = "Distribución de riqueza final",
     xlab = "Riqueza")

VaR_95 <- quantile(simulaciones, 0.05)
abline(v = VaR_95, col = "red", lwd = 2)

\end{lstlisting}
\vspace{0.4cm}



\end{frame}

%------------------------------------------------

\begin{frame}{Interpretación económica}

\textbf{Pregunta}

\vspace{0.3cm}
¿Dónde está el riesgo?

\vspace{0.4cm}

\textbf{Respuesta esperada}

\begin{itemize}
\item En la cola izquierda de la distribución
\item Donde se concentran pérdidas extremas
\item Representa escenarios adversos de mercado
\end{itemize}

\vspace{0.4cm}

\textbf{Intuición}

\begin{itemize}
\item No importa solo el promedio
\item El riesgo está en eventos poco probables pero severos
\end{itemize}

\end{frame}

%------------------------------------------------

\begin{frame}{Conexión con pensiones y stress testing}

\textbf{Aplicación previsional}

\begin{itemize}
\item Caídas en retornos afectan acumulación de fondos
\item Impacto directo en tasas de reemplazo
\end{itemize}

\vspace{0.4cm}

\textbf{Baseline vs adverso}

\begin{itemize}
\item Baseline: distribución centrada, menor dispersión
\item Adverso: mayor volatilidad, colas más pesadas
\end{itemize}

\vspace{0.4cm}

\textbf{Relevancia regulatoria}

\begin{itemize}
\item Supervisores analizan distribución completa, no solo media
\item Identificación de riesgos extremos (tail risk)
\item Base para políticas de mitigación
\end{itemize}

\end{frame}

%------------------------------------------------
\begin{frame}{Integración Final}

\textbf{Pipeline completo}

\begin{enumerate}
\item Modelo de acumulación
\item Introducción de incertidumbre
\item Simulación Monte Carlo
\item Medición de riesgo (VaR)
\item Visualización de resultados
\end{enumerate}

\vspace{0.5cm}

\textbf{Interpretación:}
Esto es un stress testing simplificado de un sistema de pensiones

\end{frame}

%------------------------------------------------
\begin{frame}{Evaluación Final}

\textbf{Responder:}

\begin{itemize}
\item ¿Qué es Monte Carlo?
\item ¿Qué representa un vector?
\item ¿Qué mide el VaR?
\end{itemize}

\vspace{0.5cm}

\centering
\textbf{Conclusión:}  
R permite transformar incertidumbre en decisiones de política previsional

\end{frame}

%------------------------------------------------

\end{document}
